% ARCHIVED at v3.92 (2026-09-24) from chapters/05_setup.tex and 06_results.tex -- author: "final parts are done and
% processed by the DA agent, use the finalize the scurve of thesis". Verbatim, not part of the build.

% ---- Ch5 fig:env-uav caption ----
  green line at the far end is the finish line, at the end of the walls on UAV-corridor and UAV-s-curve
  (UAV-s-curve is scored $0.2$\,m further on, at its goal point). The test-time constraints of UAV-corridor

% ---- Ch5 success item ----
UAV-s-curve, the tolerances with which the demonstrations of each scene were accepted. On UAV-corridor
    the finish line is the end of the corridor, $x=2.0$\,m, where the walls and both test-time constraints
    end; a flight is scored the moment it passes it and flies on as before. The line through the route's
    goal point, $0.8$\,m further on, is not used for that scene:

% ---- Ch5 finish-line srcnote ----
\srcnote{Success at the end of the corridor: \texttt{eval\_mix\_uav.py} \texttt{SCENE\_CLEAR\_LINE\_X =
    \{'corridor': 2.0\}} (author's go-ahead, 24-09; tested after each control step's physics, like the old
    latch; no new early stop; the goal-plane rule kept as \texttt{success.relaxed\_goal\_line};
    \texttt{logs\_in\_develop/Gen15/U19/CHANGELOG\_20260924\_U19\_corridor\_finish\_line\_coding2.md}).

% ---- Ch5 R44 status srcnote ----
chose per-step projection only there (option (ii), v3.80): projection R44b (job 26195) and controller
R44c (C0 job 26196, C1 after R44b) of \texttt{data\_status/PENDING\_20260923\_uav\_scurve\_R44\_raw\_first.md}
are pending.

% ---- 6.3 intro pending sentence ----
Each scene closes on its own conclusion.
Cells printed \emph{pending} are runs that have not landed (the last cell of R44c in
\texttt{data\_status/PENDING\_20260923\_uav\_scurve\_R44\_raw\_first.md}; R44b and the first R44c cell landed at
v3.91, the corridor runs, R33, at v3.80).

% ---- caveat: pilot paragraph, pilot figure, pending comparison ----
The pilot flights that opened the question are drawn in \autoref{fig:uav-scurve-paths}: the plans of
analytic average-velocity matching at $\nfe=10$, executed without projection, under each controller.
Under the cascaded geometric controller none of ten flights crosses the finish line. Each stays at
altitude until it loses attitude control and inverts, between control steps $395$ and $421$, about
$12$\,s into the flight, and ends on average $2.81$\,m from the goal point of its route. Under MuJoCo
MPC the same plans end $0.299$, $0.298$ and $0.294$\,m from that point; two of the three flights cross
the finish line, and the third reaches the finish line's $x$ with a collision. Those are three flights
against ten, and they are the pilot. Both controllers act on the same vehicle with the same nominal
gains, which no evaluation of this thesis varies; the cascaded geometric flights keep $0.29$ to $0.34$\,m
behind their commanded position, the lead of that position along the path rather than a sideways error, and
none of them advances past the second turn.

\begin{figure}[htpb]
  \centering
  \includegraphics[width=\linewidth]{fig_uav_scurve_paths}
  \caption[UAV-s-curve: the pilot flights under the two tracking controllers]{UAV-s-curve from above,
  drawn like \autoref{fig:constraints-uav}: the unprojected plans of MeanFM at $\nfe=10$ under MuJoCo
  MPC (three flights) and under the cascaded geometric controller (ten flights). Green: success; red:
  no success; a cross at the end of a path marks a flight that collided or that the divergence guard
  ended. \outdated{three flights against ten: the two panels are not comparable. The figure is replaced by
  the ten-flight cells of \autoref{tab:uav-controller} once they land (R44c).}}
  \label{fig:uav-scurve-paths}
\end{figure}
\outdated{\autoref{fig:uav-scurve-paths} compares three MuJoCo-MPC flights with ten cascaded-geometric
flights of an earlier configuration (MeanFM, $\nfe=10$); it is kept only until the ten-flight controller
comparison on FM at $\nfe=1$ lands (R44c, \autoref{tab:uav-controller}), and no conclusion rests on it.}
\dataref{DA\_in\_Paper/data/uav\_paths.json, plotting/extract/uav\_paths.py; runs
Emf\_K10\_mpc4\_mjpc\_T0.5\_u7hg (three flights, two crossing) and Emf\_K10\_mpc4\_pid\_stopgo\_T0.5\_u7hg
(ten flights, all ten aborted \texttt{inverted}). Panel titles renamed at v3.71 (\enquote{Proportional
stop-go} was a run-tag name).}

The comparison of record is made on instantaneous-velocity matching at one evaluation, the configuration of
\autoref{tab:uav-scurve-projection}, at ten flights per cell (\autoref{tab:uav-controller}): its plans,
unprojected and under per-step projection with random selection, under each controller. On the unprojected
plans MuJoCo MPC removes the one inversion and brings every flight within $0.3$\,m of the goal point, a mean
distance of $0.297$ against $0.571$\,m; nine flights of ten succeed under both controllers, the tenth under
MuJoCo MPC being scored unsafe for contact. It does not remove the violations: it follows the commanded path
less closely, $5$\,cm on average and up to $29$\,cm, the plans it is given, generated from its own states, cut
the corner deeper, and the violating steps rise from $23.0$ to $41.5$ per flight. The instability at the second
turn is therefore the controller's limit, and the violations are the plan's. Under the cascaded geometric
controller the projected plans cross on five flights of ten; the same plans under MuJoCo MPC are pending.

% ---- caveat: projected pilot paragraph ----
Changing the controller is not the whole answer once projection changes the plan. In the pilot, the
per-step-projected plans of MeanFM at $\nfe=10$ on the plain constraint set cross the finish line on
none of three flights under MuJoCo MPC and on none of ten under the cascaded geometric controller, and
no flight is free of violations under either. The unprojected pilot isolates a limit of the
controller; the projected one is a failure of the planner, the projection and the controller together.

% ---- caveat closing sentence ----
MuJoCo MPC is used here to locate the limit, and it is not adopted: on the selected plans it removes the
inversion and brings every flight within $0.3$\,m of the goal point while violating more, and in the pilot it
turns flights that never arrive into flights that arrive, at about twenty-four times the cost per control
step of the controller used everywhere else in this thesis.

% ---- s-curve conclusion (provisional, v3.91) ----
\provisional{Written on \autoref{tab:uav-scurve}, \autoref{tab:uav-scurve-projection} and three of the four cells
of \autoref{tab:uav-controller}. The last cell, the projected plans under MuJoCo MPC (R44c), and the rebuilt
\autoref{fig:uav-scurve-paths} are pending, and this conclusion is completed when they land.}

\provisional{UAV-s-curve takes the tracking controller onto an extreme trajectory, and it separates two limits.
The first is the controller's: stability at the second turn. Most flights of the average-velocity models and of
the baseline invert there under the cascaded geometric controller, just past the crossover, $11.6$ to $15.9$\,s
into the flight, and on the selected plans MuJoCo MPC removes the inversion and brings every flight within
$0.3$\,m of the goal point, at about twenty-four times the cascaded law's cost per control step --- twelve times
the cost of generating the plan, and some seven times what the operating point of D3IL-avoiding spends on an
entire action. The second limit is the plan's: every commanded path cuts the second inside corner of the
crossover, $8$ to $19$\,cm inside its keep-out, where the demonstrated route clears it by $0.121$\,m, and the
cascaded geometric controller follows it to within about a centimetre, so no flight is free of violations.
Neither remedy removes that cut as configured: MuJoCo MPC follows the command less closely and violates more,
and per-step projection, which binds the plan's measured positions as \ac{DPCC} does rather than the commanded
one, leaves the cut in place and crosses on five flights of ten instead of nine.}

\provisional{The demonstrations are simple here too, one route of straight segments joined by arcs
(\autoref{tab:uav-demos}), and on them the average-velocity objectives are not ahead at the budget they are
built for: at one evaluation instantaneous-velocity matching crosses the finish line on nine flights of ten,
the consistency-interpolated model on six and analytic average-velocity matching on three, and for every
flow-based model the crossings fall as the budget grows. Since almost every flight lost is lost to
inversion, this order is read for the controller, not for the models alone.}
\guard{Instantaneous-velocity matching is sampled from initial noise of half its training scale
(\autoref{sec:setup:protocol:avoiding}); that bears on its order against the average-velocity models, not on
the controller comparison, which flies the same plans under both controllers.}
\dataref{v3.91: R44b and C0 (\texttt{DA\_20260924\_scurve\_R44bc\_projection\_controller.md}), the verdict of its summary and \S3--\S6, \S8.3; C1 and the
rebuilt figure pending (\S8). v3.90 (author, \texttt{GUIDE\_20260924\_results\_storyline\_author.md}: \enquote{in the
s\_curve we deep discuss the controller in extreme trajs}); v3.87 (author: \enquote{claim under simple data case it
is not so as we expected \ldots\ prefill part, mark as danger sign is ok}). \autoref{tab:uav-scurve} (R44a);
inversions: $54$ of the $57$ flow-based flights lost, $48$ of them at $x$ $0.2$--$1.2$\,m; the corner cut and the
cross-track: the note to \autoref{tab:uav-scurve}; cost: \autoref{tab:uav-controller-cost} ($125.4$ against
$5.2$\,ms; planner $10.4$\,ms; $125.4/18.1 = 6.9$).}

% ---- 6.4 hole ----

\hole{UAV-s-curve joins this section when the last cell of its controller comparison (R44c) and its rebuilt
figure land; until then it is left out of the tables below.}
